%! Skills Focus: Selecting an Appropriate Inference Procedure for Categorical Data — Unit 8, Lesson 8.7 — Guided Notes (Answer Key)
\documentclass[10pt]{article}
\usepackage{apstats-article}
\usepackage{apstats-key}   % pulls in -boxes; do NOT also load -boxes

\begin{document}

\pageheader{Unit 8, Lesson 8.7}{Guided Notes --- Answer Key}
\namedateperiod

\vspace{0.10in}

\begin{objectivebox}
By the end of this lesson, I will be able to:
\begin{itemize}
  \item \ansline{Identify the correct categorical inference procedure for a scenario.}
  \item \ansline{Distinguish between $z$-procedures and chi-square procedures.}
  \item \ansline{State $H_0$ in the right language for each categorical procedure.}
\end{itemize}
\end{objectivebox}

\begin{vocabbox}
\begin{description}
  \item[\termblanklong{Categorical inference procedure}]
        Any inference method applied to \makebox[4cm][l]{\ans{categorical}} data. The first diagnostic question is:
        does the outcome have \makebox[2cm][l]{\ans{two}} categories or \makebox[3cm][l]{\ans{three or more}}?

  \item[\termblanklong{Chi-square goodness-of-fit (GOF) test}]
        Used when there is \makebox[2cm][l]{\ans{one}} sample and \makebox[2cm][l]{\ans{three}} or more categories. Tests
        whether the observed distribution matches a \makebox[5cm][l]{\ans{claimed}} distribution.

  \item[\termblanklong{Chi-square test for homogeneity}]
        Used when there are \makebox[2.5cm][l]{\ans{two}} or more \makebox[3cm][l]{\ans{independent}} samples and one categorical
        variable. Tests whether the \makebox[4.5cm][l]{\ans{variable's distribution}} is the same across groups.

  \item[\termblanklong{Chi-square test for independence}]
        Used when there is \makebox[2cm][l]{\ans{one}} sample and \makebox[2cm][l]{\ans{two}} categorical variables recorded
        for each subject. Tests whether there is an \makebox[4cm][l]{\ans{association}} between the variables.
\end{description}
\end{vocabbox}

\begin{hookbox}
A researcher surveys students about their transportation to school (Walk, Bus, Car, Bike).
\begin{enumerate}
  \item \textit{Scenario A:} She surveys students at one school and wants to test whether
        walking, bus, car, and biking are equally common.\\
        What procedure? \makebox[8cm][l]{\ans{chi-square goodness-of-fit test}}
  \item \textit{Scenario B:} She surveys students at two schools (separate random samples)
        and wants to test whether the distribution of transportation modes is the same at both.\\
        What procedure? \makebox[8cm][l]{\ans{chi-square test for homogeneity}}
  \item \textit{Scenario C:} She takes one random sample and records both transportation mode
        and whether each student arrives on time (yes/no). She tests for association.\\
        What procedure? \makebox[8cm][l]{\ans{chi-square test for independence}}
\end{enumerate}
\end{hookbox}

\begin{notesbox}{Categorical Inference Decision Table}
\textbf{Fill in the procedure name and key condition:}

\vspace{0.06in}
\renewcommand{\arraystretch}{2.2}
\begin{tabularx}{\linewidth}{@{} p{0.20\linewidth} p{0.20\linewidth} p{0.12\linewidth} p{0.20\linewidth} X @{}}
  \rowcolor{sky}
  \textbf{Data Structure} & \textbf{Number of Categories} & \textbf{Goal} & \textbf{Procedure Name} & \textbf{Key Condition} \\
  \hline
  One proportion, one sample & Two (binary) & CI & \ans{1-sample $z$-interval for $p$} & \ans{Large Counts: $n\hat p, n(1{-}\hat p) \ge 10$} \\
  One proportion, one sample & Two (binary) & Test & \ans{1-sample $z$-test for $p$} & \ans{Large Counts: $np_0, n(1{-}p_0) \ge 10$} \\
  Difference in proportions, two samples & Two (binary) & CI & \ans{2-sample $z$-interval for $p_1{-}p_2$} & \ans{Large Counts: all four $\ge 10$} \\
  Difference in proportions, two samples & Two (binary) & Test & \ans{2-sample $z$-test for $p_1{-}p_2$} & \ans{Large Counts (pooled) $\ge 10$} \\
  One sample, claimed distribution & Three or more & Test & \ans{$\chi^2$ goodness-of-fit test} & \ans{All expected counts $\ge 5$} \\
  Two or more groups (sep.\ samples), one variable & Three or more & Test & \ans{$\chi^2$ test for homogeneity} & \ans{All expected counts $\ge 5$} \\
  One sample, two variables & Any & Test & \ans{$\chi^2$ test for independence} & \ans{All expected counts $\ge 5$} \\
\end{tabularx}
\end{notesbox}

\begin{notesbox}{Homogeneity vs.\ Independence --- How to Tell}
\renewcommand{\arraystretch}{1.8}
\begin{tabularx}{\linewidth}{@{} p{0.22\linewidth} X X @{}}
  \rowcolor{sky}
  & \textbf{Homogeneity} & \textbf{Independence} \\ \hline
  \textbf{Sampling design} & \ans{separate random samples} & \ans{one random sample} \\
  \textbf{Number of variables} & \makebox[3cm][l]{\ans{one}} & \makebox[3cm][l]{\ans{two}} \\
  \textbf{$H_0$ language} & ``The \makebox[2.5cm][l]{\ans{distribution}} is the same across all \makebox[2cm][l]{\ans{groups}}.'' & ``There is no \makebox[3cm][l]{\ans{association}} between [var 1] and [var 2].'' \\
  \textbf{Same $\chi^2$ statistic?} & \makebox[2cm][l]{\ans{yes}} & \makebox[2cm][l]{\ans{yes}} \\
\end{tabularx}
\end{notesbox}

\begin{practicebox}
\textbf{Partner Practice.} \;
A social scientist randomly selects 50 adults aged 18--34 and 50 adults aged 35--54. He asks
each person to rate their social media use as Low, Moderate, or High. He wants to test whether
the distribution of social media use is the same for both age groups.

\begin{enumerate}
  \item Binary or multi-category outcome? \makebox[4cm][l]{\ans{multi-category}} \quad Number of groups: \makebox[2cm][l]{\ans{2}}
  \item Procedure name: \makebox[9cm][l]{\ans{chi-square test for homogeneity}}
  \item Form of $H_0$: \ansline{Use is distributed the same for ages 18--34 and 35--54.}
  \item Key condition to check (beyond Random and 10\%): \ansline{All expected counts $\ge 5$.}
\end{enumerate}
\end{practicebox}

\end{document}
